
% ────────────────────────────────────────────────────────
\section{Aplicações e Comprimento do Arco}

% ────────────────────────────────────────────────────────
\subsection{Enquadramento Teórico}

\subsubsection*{Área entre Curvas}

Se $f(x)\geq g(x)$ em $[a,b]$, a área da região delimitada é:
\[
  A = \int_a^b \bigl[f(x)-g(x)\bigr]\,\dd x.
\]

\subsubsection*{Volume de um Sólido de Revolução (Método dos Discos)}

Ao rodar $y=f(x)\geq 0$ em torno do eixo $Ox$ de $a$ a $b$:
\[
  V = \pi\int_a^b [f(x)]^2\,\dd x.
\]

\subsubsection*{Comprimento do Arco}

O comprimento da curva $y=f(x)$ em $[a,b]$:
\[
  L = \int_a^b \sqrt{1+[f'(x)]^2}\,\dd x.
\]

Para uma curva paramétrica $(x(t),y(t))$, $t\in[\alpha,\beta]$:
\[
  L = \int_\alpha^\beta \sqrt{[\dot x(t)]^2+[\dot y(t)]^2}\,\dd t.
\]

% ────────────────────────────────────────────────────────
\subsection{Exemplos Resolvidos}

\begin{exbox}{6 --- Área entre Curvas}
  Determine a área entre $f(x)=x^2$ e $g(x)=x$ em $[0,1]$.
\end{exbox}
\begin{solbox}
  Como $x\geq x^2$ em $[0,1]$:
  \[
    A = \int_0^1(x-x^2)\,\dd x = \left[\frac{x^2}{2}-\frac{x^3}{3}\right]_0^1
      = \frac{1}{2}-\frac{1}{3} = \frac{1}{6}.
  \]
\end{solbox}

\begin{exbox}{7 --- Volume de uma Esfera}
  Volume da esfera de raio $r$: roda-se $f(x)=\sqrt{r^2-x^2}$ em
  $[-r,r]$ em torno do eixo $Ox$.
\end{exbox}
\begin{solbox}
  \textbf{Passo 1 -- Identificar a função.} A semicircunferência superior de
  raio $r$ é dada por $f(x)=\sqrt{r^2-x^2}$, com $x\in[-r,r]$. Ao rodar
  este arco em torno do eixo $Ox$ obtém-se uma esfera completa.

  \textbf{Passo 2 -- Aplicar a fórmula dos discos.}
  \[
    V = \pi\int_{-r}^r [f(x)]^2\,\dd x = \pi\int_{-r}^r (r^2-x^2)\,\dd x.
  \]

  \textbf{Passo 3 -- Calcular o integral.}
  \[
    V = \pi\left[r^2 x - \frac{x^3}{3}\right]_{-r}^r
      = \pi\left[\left(r^3-\frac{r^3}{3}\right)-\left(-r^3+\frac{r^3}{3}\right)\right]
      = \pi\cdot\frac{4r^3}{3} = \frac{4}{3}\pi r^3.
  \]
\end{solbox}

\begin{exbox}{8 --- Comprimento do Arco}
  Determine o comprimento de $y=\dfrac{2}{3}x^{3/2}$ de $x=0$ a $x=3$.
\end{exbox}
\begin{solbox}
  $f'(x)=x^{1/2}$, logo $1+[f']^2=1+x$.
  \[
    L = \int_0^3\sqrt{1+x}\,\dd x = \left[\tfrac{2}{3}(1+x)^{3/2}\right]_0^3
      = \tfrac{2}{3}(8-1) = \tfrac{14}{3}.
  \]
\end{solbox}

% ────────────────────────────────────────────────────────
\subsection{Exercícios Práticos}

\begin{exercisebox}{1}
  Calcule os seguintes integrais indefinidos:
  \begin{enumerate}[label=(\alph*)]
    \item $\displaystyle\int (3x^2-2x+5)\,\dd x$
    \item $\displaystyle\int \frac{x}{x^2+4}\,\dd x$
    \item $\displaystyle\int x^2 e^x\,\dd x$
    \item $\displaystyle\int \frac{x+1}{x^2+x-2}\,\dd x$
    \item $\displaystyle\int \cos^2 x\,\dd x$
  \end{enumerate}
\end{exercisebox}
\begin{solbox}
  \textbf{(a)} Aplica-se a linearidade e a regra da potência:
  \[
    \int (3x^2-2x+5)\,\dd x = 3\cdot\frac{x^3}{3} - 2\cdot\frac{x^2}{2} + 5x + C
    = x^3 - x^2 + 5x + C.
  \]

  \textbf{(b)} Reconhece-se a forma $\displaystyle\int \frac{f'(x)}{f(x)}\,\dd x$.
  Seja $u=x^2+4$, $du=2x\,\dd x$, logo $x\,\dd x=\dfrac{du}{2}$:
  \[
    \int \frac{x}{x^2+4}\,\dd x = \frac{1}{2}\int\frac{du}{u}
    = \frac{1}{2}\ln|u|+C = \frac{1}{2}\ln(x^2+4)+C.
  \]

  \textbf{(c)} Aplica-se primitivação por partes duas vezes (LIATE: $u=x^2$,
  $dv=e^x\,\dd x$):
  \begin{align*}
    &\textbf{1.ª vez:}\quad u=x^2,\ dv=e^x\,\dd x \Rightarrow du=2x\,\dd x,\ v=e^x.\\
    &\int x^2 e^x\,\dd x = x^2 e^x - 2\int x e^x\,\dd x.\\
    &\textbf{2.ª vez:}\quad u=x,\ dv=e^x\,\dd x \Rightarrow du=\dd x,\ v=e^x.\\
    &\int x e^x\,\dd x = xe^x - e^x + C.\\
    &\textbf{Resultado:}\quad \int x^2 e^x\,\dd x = x^2 e^x - 2(xe^x - e^x)+C
      = e^x(x^2-2x+2)+C.
  \end{align*}

  \textbf{(d)} Fatoriza-se o denominador: $x^2+x-2=(x+2)(x-1)$.
  Decomposição em frações parciais:
  \[
    \frac{x+1}{(x+2)(x-1)} = \frac{A}{x+2}+\frac{B}{x-1}.
  \]
  Multiplicando ambos os membros por $(x+2)(x-1)$: $x+1 = A(x-1)+B(x+2)$.
  \begin{itemize}
    \item $x=1$: $2=3B \Rightarrow B=\tfrac{2}{3}$.
    \item $x=-2$: $-1=-3A \Rightarrow A=\tfrac{1}{3}$.
  \end{itemize}
  \[
    \int \frac{x+1}{x^2+x-2}\,\dd x = \frac{1}{3}\ln|x+2|+\frac{2}{3}\ln|x-1|+C.
  \]

  \textbf{(e)} Usa-se a identidade do ângulo duplo: $\cos^2 x = \dfrac{1+\cos 2x}{2}$.
  \[
    \int \cos^2 x\,\dd x = \int \frac{1+\cos 2x}{2}\,\dd x
    = \frac{x}{2} + \frac{\sin 2x}{4} + C.
  \]
\end{solbox}

\begin{exercisebox}{2}
  Calcule os seguintes integrais definidos:
  \begin{enumerate}[label=(\alph*)]
    \item $\displaystyle\int_1^e \ln x\,\dd x$
    \item $\displaystyle\int_0^1 x\sqrt{1-x^2}\,\dd x$
    \item $\displaystyle\int_0^{\pi/2}\sin^2 x\,\dd x$
  \end{enumerate}
\end{exercisebox}
\begin{solbox}
  \textbf{(a)} Primitivação por partes: $u=\ln x$, $dv=\dd x$, logo
  $du=\dfrac{1}{x}\,\dd x$, $v=x$.
  \[
    \int_1^e \ln x\,\dd x = \Big[x\ln x\Big]_1^e - \int_1^e x\cdot\frac{1}{x}\,\dd x
    = \Big[x\ln x\Big]_1^e - \Big[x\Big]_1^e
    = (e\cdot 1 - 1\cdot 0)-(e-1) = 1.
  \]

  \textbf{(b)} Substituição: $u=1-x^2$, $du=-2x\,\dd x$.
  Limites: $x=0\Rightarrow u=1$; $x=1\Rightarrow u=0$.
  \[
    \int_0^1 x\sqrt{1-x^2}\,\dd x = -\frac{1}{2}\int_1^0 \sqrt{u}\,\dd u
    = \frac{1}{2}\int_0^1 u^{1/2}\,\dd u
    = \frac{1}{2}\left[\frac{2}{3}u^{3/2}\right]_0^1 = \frac{1}{3}.
  \]

  \textbf{(c)} Identidade: $\sin^2 x = \dfrac{1-\cos 2x}{2}$.
  \[
    \int_0^{\pi/2}\sin^2 x\,\dd x = \int_0^{\pi/2}\frac{1-\cos 2x}{2}\,\dd x
    = \left[\frac{x}{2}-\frac{\sin 2x}{4}\right]_0^{\pi/2}
    = \frac{\pi}{4}.
  \]
\end{solbox}

\begin{exercisebox}{3}
  Determine a área da região delimitada por $y=x^3$ e $y=x$.
\end{exercisebox}
\begin{solbox}
  \textbf{Passo 1 -- Pontos de interseção.} $x^3=x \Rightarrow x(x^2-1)=0
  \Rightarrow x\in\{-1,0,1\}$.

  \textbf{Passo 2 -- Sinal da diferença.} Em $(-1,0)$: $x^3-x = x(x^2-1)$.
  Teste em $x=-\tfrac{1}{2}$: $(-\tfrac{1}{2})({\tfrac{1}{4}-1})>0$,
  logo $x^3\geq x$ em $(-1,0)$.
  Em $(0,1)$: teste em $x=\tfrac{1}{2}$: resultado negativo, logo $x\geq x^3$.

  \textbf{Passo 3 -- Cálculo por simetria.} A função $x^3-x$ é ímpar, pelo que
  a área em $(-1,0)$ iguala a área em $(0,1)$:
  \[
    A = 2\int_0^1(x-x^3)\,\dd x = 2\left[\frac{x^2}{2}-\frac{x^4}{4}\right]_0^1
    = 2\cdot\frac{1}{4} = \frac{1}{2}.
  \]
\end{solbox}

\begin{exercisebox}{4}
  Determine o volume do sólido obtido ao rodar a região delimitada por
  $y=\sqrt{x}$, $x=4$ e $y=0$ em torno do eixo $Ox$.
\end{exercisebox}
\begin{solbox}
  \textbf{Passo 1 -- Identificar o método.} A região é delimitada por baixo por
  $y=0$ e por cima por $y=\sqrt{x}$, com $x\in[0,4]$. Aplica-se o método dos
  discos.

  \textbf{Passo 2 -- Aplicar a fórmula.}
  \[
    V = \pi\int_0^4 [\sqrt{x}]^2\,\dd x = \pi\int_0^4 x\,\dd x.
  \]

  \textbf{Passo 3 -- Calcular.}
  \[
    V = \pi\left[\frac{x^2}{2}\right]_0^4 = \pi\cdot\frac{16}{2} = 8\pi.
  \]
\end{solbox}

\begin{exercisebox}{5}
  Calcule o comprimento de arco da curva $y = \ln(\cos x)$ de $x=0$ a
  $x=\pi/4$.
\end{exercisebox}
\begin{solbox}
  \textbf{Passo 1 -- Derivada.}
  \[
    y' = \frac{-\sin x}{\cos x} = -\tan x.
  \]

  \textbf{Passo 2 -- Expressão sob a raiz.}
  \[
    1+[y']^2 = 1+\tan^2 x = \sec^2 x.
  \]

  \textbf{Passo 3 -- Integral do comprimento.}
  \[
    L = \int_0^{\pi/4}\sqrt{\sec^2 x}\,\dd x = \int_0^{\pi/4}|\sec x|\,\dd x
    = \int_0^{\pi/4}\sec x\,\dd x.
  \]
  (Note-se que $\sec x > 0$ em $[0,\pi/4]$.)

  \textbf{Passo 4 -- Calcular $\displaystyle\int\sec x\,\dd x$.}
  Usando a primitiva conhecida $\displaystyle\int\sec x\,\dd x = \ln|\sec x+\tan x|+C$:
  \[
    L = \Big[\ln|\sec x + \tan x|\Big]_0^{\pi/4}
    = \ln\!\left(\sqrt{2}+1\right) - \ln(1+0)
    = \ln(\sqrt{2}+1).
  \]
\end{solbox}

% ────────────────────────────────────────────────────────
\subsection{Questões Propostas para Defesa Oral}

\begin{oralbox}
  \textbf{Q1.} Enuncie e demonstre o Teorema Fundamental do Cálculo Integral
  (Parte 2). Qual o papel da continuidade?
\end{oralbox}

\begin{oralbox}
  \textbf{Q2.} Explique a diferença entre integral indefinido e integral definido.
  Por que razão a constante arbitrária $C$ desaparece no caso definido?
\end{oralbox}

\begin{oralbox}
  \textbf{Q3.} Quando é que a primitivação por partes é mais útil? Dê um exemplo
  em que seja necessário aplicá-la duas vezes.
\end{oralbox}

\begin{oralbox}
  \textbf{Q4.} Descreva geometricamente o que representa o integral definido
  $\displaystyle\int_a^b f(x)\,\dd x$ quando $f$ toma valores negativos em parte de $[a,b]$.
\end{oralbox}

\begin{oralbox}
  \textbf{Q5.} Deduza a fórmula do comprimento de arco a partir de primeiros
  princípios, recorrendo ao Teorema de Pitágoras.
\end{oralbox}

\begin{oralbox}
  \textbf{Q6.} Compare o método dos discos e o método das cascas para o cálculo
  de volumes. Em que circunstâncias se prefere cada um?
\end{oralbox}

\begin{oralbox}
  \textbf{Q7.} Qual a relação entre a área sob uma curva e uma soma de Riemann?
  Descreva a convergência.
\end{oralbox}